% Chunk: Chapter 5 back matter---Problems and References.
% Source: Weinberg QFT Vol. I, physical PDF pages 278--281
%         (printed pages 255--258), beginning at the Problems heading.
% Status: source-reviewed, compile-clean, visually verified, and frozen.
% Problems: 7. References: 18. Figures: none. Uncertainties: none.

\chapterexercisehook{05}

\chapterbackmatter{References}

\begin{enumerate}
\item\label{ref:5.1} The point of view adopted in this chapter was presented in a series of
papers: S. Weinberg, \emph{Phys. Rev.} \textbf{133}, B1318 (1964);
\textbf{134}, B882 (1964); \textbf{138}, B988 (1965); \textbf{181}, 1893
(1969). A similar approach has been followed in unpublished lectures by
E. Wichmann.

\item\label{ref:5.2} N. Bohr and L. Rosenfeld, \emph{Kgl. Danske Vidensk. Selskab Mat.-Fys.
Medd.}, No. 12 (1933) (translation in \emph{Selected Papers of Leon
Rosenfeld}, ed. by R. S. Cohen and J. Stachel (Reidel, Dordrecht, 1979));
\emph{Phys. Rev.} \textbf{78}, 794 (1950).

\item\label{ref:5.3} P. A. M. Dirac, \emph{Proc. Roy. Soc. (London)} A\textbf{117}, 610
(1928).

\item\label{ref:5.4} E. Cartan, \emph{Bull. Soc. Math. France} \textbf{41}, 53 (1913).

\item\label{ref:5.5} See, e.g., J. M. Jauch and F. Rohrlich, \emph{The Theory of Photons and
Electrons} (Addison-Wesley, Cambridge, MA, 1955): Appendix A2.

\item\label{ref:5.6} See, e.g., H. Georgi, \emph{Lie Algebras in Particle Physics}
(Benjamin/Cummings, Reading, MA, 1982): pp. 15, 198. The original reference
is I. Schur, \emph{Sitz. Preuss. Akad.}, p. 406 (1905).

\item\label{ref:5.7} T. D. Lee and C. N. Yang, \emph{Phys. Rev.} \textbf{104}, 254 (1956).

\item\label{ref:5.8} See, e.g., B. L. van der Waerden, \emph{Die gruppentheoretische Methode
in der Quantenmechanik} (Springer Verlag, Berlin, 1932); G. Ya. Lyubarski,
\emph{The Applications of Group Theory in Physics}, translated by S. Dedijer
(Pergamon Press, New York, 1960).

\item\label{ref:5.9} W. Rarita and J. Schwinger, \emph{Phys. Rev.} \textbf{60}, 61 (1941).

\item\label{ref:5.10} See, e.g., A. R. Edmonds, \emph{Angular Momentum in Quantum Mechanics}
(Princeton University Press, Princeton, 1957): Chapter 3.

\item\label{ref:5.11} S. Weinberg, \emph{Phys. Rev.} \textbf{181}, 1893 (1969), Section V.

\item\label{ref:5.12} M. Fierz, \emph{Helv. Phys. Acta} \textbf{12}, 3 (1939); W. Pauli,
\emph{Phys. Rev.} \textbf{58}, 716 (1940). Non-perturbative proofs in
axiomatic field theory were given by G. Lüders and B. Zumino,
\emph{Phys. Rev.} \textbf{110}, 1450 (1958) and N. Burgoyne,
\emph{Nuovo Cimento} \textbf{8}, 807 (1958). Also see R. F. Streater and
A. S. Wightman, \emph{PCT, Spin \& Statistics, and All That} (Benjamin,
New York, 1968).

\item\label{ref:5.13} Fields in the $(j,0)+(0,j)$ representation were introduced by H. Joos,
\emph{Fortschr. Phys.} \textbf{10}, 65 (1962); S. Weinberg,
\emph{Phys. Rev.} \textbf{133}, B1318 (1964).

\item\label{ref:5.14} A. R. Edmonds, Ref.~\hyperref[ref:5.10]{10}, or M. E. Rose, \emph{Elementary Theory of
Angular Momentum} (John Wiley \& Sons, New York, 1957): Chapter III.

\item\label{ref:5.15} G. Velo and D. Zwanziger, \emph{Phys. Rev.} \textbf{186}, 1337 (1969);
\textbf{188}, 2218 (1969); A. S. Wightman, in \emph{Proceedings of the Fifth
Coral Gables Conference on Symmetry Principles at High Energy}, ed. by
T. Gudehus, G. Kaiser, and A. Perlmutter (Gordon and Breach, New York,
1969); B. Schroer, R. Seiler, and J. A. Swieca, \emph{Phys. Rev. D}
\textbf{2}, 2927 (1970); and other references quoted therein.

\item\label{ref:5.16} C. R. Nappi and L. Witten, \emph{Phys. Rev. D} \textbf{40}, 1095
(1989); P. C. Argyres and C. R. Nappi, \emph{Phys. Lett.} B\textbf{224}, 89
(1989).

\item\label{ref:5.17} G. Lüders, \emph{Kong. Dansk. Vid. Selskab, Mat.-Fys. Medd.}
\textbf{28}, 5 (1954); \emph{Ann. Phys.} \textbf{2}, 1 (1957); W. Pauli,
\emph{Nuovo Cimento} \textbf{6}, 204 (1957). When Lüders first considered
how the inversions are related, it was still taken for granted that $P$ is
conserved, so his theorem stated that $C$ conservation is equivalent to $T$
invariance.

\item\label{ref:5.18} R. Jost, \emph{Helv. Phys. Acta} \textbf{30}, 409 (1957); F. J. Dyson,
\emph{Phys. Rev.} \textbf{110}, 579 (1958). Also see Streater and Wightman,
Ref.~\hyperref[ref:5.12]{12}.
\end{enumerate}
